\chapter*{常见缩写和术语先导}
\addcontentsline{toc}{chapter}{常见缩写和术语先导}

第三册的缩写和工程词最密，因为它同时涉及模型、编译器、CPU/GPU 后端和服务运行时。下面这页不是让你背名词，而是把字母、形状词和工程对象先和本书主线对齐。

\begin{description}
  \item[CPU/GPU/ISA/API/ABI/SDK/GDB] CPU 是中央处理器；GPU 是图形处理器或通用并行加速器；ISA 是指令集架构；API 是源码接口；ABI 是二进制接口；SDK 是软件开发工具包；GDB 是 GNU 调试器。第三册会不断把模型算子落到这些机器和工具边界。
  \item[C++20] C++ 标准版本名。本册的 descriptor、RAII、span、oracle、runtime plan 等工程示例默认以 C++20 为实现边界。
  \item[SIMD/AVX/AVX2/AVX-512] SIMD 是一条指令处理多个数据；AVX/AVX2 是 x86 向量扩展；AVX-512 提供更宽寄存器和 mask 机制。后端选择要看 shape、tail、频率和平台覆盖。
  \item[FMA/VNNI/AMX/PMU/IPC] FMA 是乘加指令；VNNI 是低精度 dot-product 相关指令能力；AMX 是 x86 tile 级矩阵扩展；PMU 提供硬件计数器；IPC 是每周期指令数。
  \item[TLB/L2/L3/LLC/HBM/DRAM/NUMA] TLB 是地址翻译缓存；L2/L3/LLC 是缓存层级；HBM 是高带宽显存；DRAM 是主存；NUMA 表示不同内存节点访问成本不同。权重流、KV cache 和 activation arena 都会碰到这些成本。
  \item[JSON/CSV/CI/I/O] JSON 常保存模型配置、tokenizer 或 manifest；CSV 常保存 benchmark 表格；CI 是持续集成；I/O 是输入输出。模型加载和性能验收离不开这些工程格式。
  \item[Unicode/UTF-8] Unicode 是统一字符集合；UTF-8 是它的常见变长字节编码。Tokenizer、prompt、stop string 和 golden case 都必须保存清楚的字节边界。
  \item[KB/KiB/MB/MiB/GB/GiB] 容量单位。KB/MB/GB 常按十进制理解，KiB/MiB/GiB 明确按 1024 进位；模型权重、KV cache、arena 和显存预算必须写清单位口径。
  \item[mmap/perf/objdump] mmap 把模型文件或匿名页映射到进程虚拟地址空间；perf 用采样和 PMU 事件看瓶颈；objdump 用来确认后端是否生成预期指令。它们分别服务加载、测量和机器码验证。
  \item[LLM] Large Language Model，大语言模型。本书关注它在本地推理时如何变成张量计算、KV cache 和请求调度。
  \item[KV cache] Key/Value cache，Transformer 解码时保存历史 token 的 K/V 张量，避免每生成一个 token 都重算全部历史。
  \item[QKV] Query、Key、Value。attention 会用 Q 查历史 K，并把权重作用到 V 上得到上下文表示。
  \item[BM/BN/BK] GEMM tile 记号。BM 表示一次处理多少 batch 行，BN 表示多少输出通道，BK 表示沿 K 维一次推进多少元素。
  \item[TILE\_N/TILE\_K] 后端实验里的 tile 常量。TILE\_N 表示 N 方向一次处理的列块，TILE\_K 表示归约维度块；它们是 schedule 参数，不是新模型结构。
  \item[QH/KVH/KVA] QH 是 query head 数；KVH 或 KVA 是 key/value head 数。它们不是固定算法名，而是 shape verifier 里必须检查的一组形状变量。
  \item[MLP/BLAS] MLP 是 Multi-Layer Perceptron，多层感知机；BLAS 是 Basic Linear Algebra Subprograms，基础线性代数接口族。Transformer 块里的 MLP 通常指 attention 后面的前馈网络。
  \item[GEMM/GEMV/FLOP/MAC] GEMM 是矩阵乘矩阵，GEMV 是矩阵乘向量；FLOP 是浮点运算次数；MAC 是 multiply-accumulate，乘加次数。prefill 更像 GEMM，decode 中很多投影更像 GEMV。
  \item[FP16/BF16/TF32/INT8/MFLOP/GFLOP/TFLOP] FP16、BF16、TF32、INT8 是常见低精度 dtype；MFLOP/GFLOP/TFLOP 是百万、十亿、万亿级浮点运算量或吞吐单位。报告必须写清 dtype 和 FLOP 口径。
  \item[RMSNorm/RoPE/SwiGLU] RMSNorm 是均方根归一化；RoPE 是旋转位置编码；SwiGLU 是一种门控前馈网络激活结构。
  \item[MHA/GQA/MQA] MHA 是多头注意力；GQA 让多个 query head 共享一组 KV head；MQA 让所有 query head 共享一组 KV head。
  \item[QK/PV] attention 中的两个主要矩阵乘。QK 计算 query 和 key 的相似度分数；PV 把 softmax 后的权重作用到 value 上。
  \item[FlashAttention] 一类分块 attention 思路。数学结果仍是 attention，但通过在线 softmax 和 tile 复用减少中间矩阵写回。
  \item[TTFT/TPOT] TTFT 是 Time To First Token，首 token 延迟；TPOT 是 Time Per Output Token，之后每个输出 token 的平均时间。
  \item[p50/p95/p99] 延迟分位数。p50 表示典型请求，p95/p99 表示尾部请求；模型服务要同时看 TTFT/TPOT 的分位数，而不能只看平均值。
  \item[SLO] Service Level Objective，服务等级目标，例如 TTFT、TPOT 或 p95 延迟阈值。
  \item[FIFO/FAQ/DDIA/RAG/BM25/HTML] FIFO 是先进先出调度；FAQ 是常见问题；DDIA 常指《Designing Data-Intensive Applications》这类数据密集系统材料；RAG 是检索增强生成；BM25 是经典稀疏关键词检索打分；HTML 是网页标记语言。
  \item[GGUF] llama.cpp 生态常见模型文件格式，通常把权重、量化信息和模型元数据打包在一个文件中。
  \item[ggml/llama.cpp/vLLM/oneDNN] ggml 和 llama.cpp 是本地 LLM 推理生态中的 C/C++ 项目；vLLM 是高吞吐 LLM 服务系统；oneDNN 是 CPU 深度学习算子库。本册提到它们时，是把工业实现当作 shape、layout、kernel 和 runtime 的对照样本。
  \item[ONNX] Open Neural Network Exchange，开放神经网络交换格式。它用于描述模型图、算子和权重，导入时必须经过 shape、dtype 和算子语义校验。
  \item[BPE/BOS/EOS/PAD/UNK/OOV] BPE 是 Byte Pair Encoding，一类 tokenizer；BOS 是句首 token，EOS 是句尾 token，PAD 是填充 token；UNK 是未知 token；OOV 是 out-of-vocabulary，词表外输入。
  \item[组合缩写] \code{BOS/EOS}、\code{GEMM/GEMV}、\code{MLIR/TVM}、\code{VNNI/AMX} 这类斜杠组合表示“把几个相关概念并列讨论”。斜杠不是新术语，读的时候拆回每个条目即可。
  \item[MoE/LoRA] MoE 是 Mixture of Experts，专家混合结构；LoRA 是 Low-Rank Adaptation，用低秩增量适配模型权重。
  \item[GPTQ/AWQ/NF4] 常见大模型权重量化名词。GPTQ 偏二阶近似逐层量化，AWQ 强调保护重要 activation 通道，NF4 是非均匀 4-bit codebook 量化格式之一。
  \item[LCQI] 本书项目里的 Local CPU Quantized Inference，指本地 CPU 量化推理实验线。它是教学工程代号，不是业界标准缩写。
  \item[LCQI\_TOKENIZER\_V1] 本册教学工程中的 tokenizer manifest 版本名，表示“LCQI 项目的第 1 版 tokenizer 合同”。它不是外部标准。
  \item[IR/SSA/pass/lowering/UB] IR 是中间表示；SSA 是静态单赋值形式；pass 是读取并分析或改写 IR 的步骤；lowering 是把高层表示逐步降到后端可执行形式；UB 是 undefined behavior，C++ 未定义行为。
  \item[AST] Abstract Syntax Tree，抽象语法树。第三册提到它时，通常是为了把传统编译器概念和 AI graph IR 做类比。
  \item[MLIR/TVM/XLA] 三类工业编译器或编译框架。本书提到它们时，重点是比较 IR、pass、schedule 和 lowering 的工程边界。
  \item[CUDA/HIP/NCCL/PCIe] CUDA 是 NVIDIA GPU 编程接口；HIP 常用于 AMD GPU 生态；NCCL 是 NVIDIA Collective Communications Library，常用于多 GPU collective 通信；PCIe 是主机和设备之间常见互连。
  \item[ACL] Arm Compute Library，Arm 生态常见 CPU/GPU 计算库。本书提到它时，重点是把算子合同映射到已有后端库。
  \item[SM/CU] GPU 的基本执行资源组。NVIDIA 常称 SM，AMD 常称 CU。
  \item[HBM/GDDR] GPU 常见显存类型。HBM 带宽高、封装近；GDDR 更常见于消费级 GPU。后端报告要区分显存带宽和主机 DRAM 带宽。
  \item[OOM/RNG/NaN] OOM 是内存不足；RNG 是随机数生成器；NaN 是 not-a-number，常提示数值计算已经失控。
  \item[KL/MSE/MMR/MRR/PII] KL divergence 和 mean squared error 是量化校准中常见指标；MMR 是最大边际相关性，用来在检索中平衡相关性和多样性；MRR 是平均倒数排名；PII 是 personally identifiable information，个人可识别信息。
  \item[UI/PR] UI 是用户界面；PR 是 pull request，代码合并请求。验收章节提到它们时，关注状态机、回归门禁和可审查证据。
  \item[KB/MB/GB] 存储容量单位，分别表示千字节、兆字节、十亿字节量级。模型大小、KV cache 和 arena 报告必须写清单位口径。
  \item[token/tokenizer/prompt/vocab] token 是模型处理的离散编号或片段；tokenizer 把文本字节映射成 token 序列；prompt 是输入给模型的上下文；vocab 是可用 token 集合。不要把用户眼里的一个汉字、一个字节和一个 token 混为一谈。
  \item[shape/dtype/layout/stride] shape 是张量各维长度；dtype 是元素类型；layout 是物理排列；stride 是相邻下标在内存中的步幅。后端优化前必须先证明这四件事解释一致。
  \item[descriptor/metadata/schema/manifest] descriptor 是结构化描述对象；metadata 是描述数据的数据；schema 是格式规则；manifest 是导入后保存模型、tokenizer、权重、量化和后端事实的机器可读清单。
  \item[reference/oracle/golden/fixture] reference 是定义正确结果的清晰实现；oracle 是判定器；golden 是已确认的期望结果；fixture 是测试输入和环境搭好的样本。
  \item[backend/runtime/kernel/microkernel] backend 是具体硬件实现层；runtime 是执行计划、内存和调度层；kernel 是某个算子的具体热路径实现；microkernel 是更小的 tile 级核心计算单元。
  \item[arena/workspace/packed weight] arena 是预先规划的大块内存；workspace 是临时工作区；packed weight 是为后端读取顺序重排后的权重。它们决定热路径是否反复分配、拷贝和解包。
  \item[prefill/decode/logits/sampling] prefill 是处理整段 prompt、建立 KV cache 的阶段；decode 是逐 token 生成阶段；logits 是下一个 token 的打分；sampling 是按 greedy、top-k、top-p 或温度等策略选 token。
  \item[tail/mask/lane/tile/block] tail 是不能整除 tile、SIMD 宽度或 group size 的尾部；mask 标记有效 lane；lane 是向量槽；tile 和 block 是分块计算单位。低比特推理的大量 bug 都出在这些边界。
  \item[stream/event/callback/backpressure] stream 是 GPU 或服务运行时的异步执行队列；event 是可记录和等待的完成点；callback 是完成后回调；backpressure 是过载时让上游减速、排队或失败。
  \item[baseline/benchmark/profile/trace] baseline 是基线；benchmark 是受控性能测量；profile 是运行画像；trace 是事件流水。第三册的优化结论必须绑定这四类证据。
  \item[ReLU/ReLU6/activation] ReLU 是把负数截到 0 的激活函数；ReLU6 进一步把结果截到 6；activation 是中间激活张量或激活函数。量化和融合时必须说明 clamp 发生在哪里。
  \item[SentencePiece/WordPiece/byte fallback] SentencePiece 和 WordPiece 是常见 tokenizer 算法族；byte fallback 是词表不覆盖时退回到字节级表示。它们决定同一段 UTF-8 文本会变成哪些 token。
  \item[LangChain/LangGraph/LlamaIndex] 常见上层 LLM 应用编排框架。它们不等于推理引擎，但会改变 prompt、tool、retrieval、stream、checkpoint 和 trace 合同。
  \item[OpenAI-style/URL/endpoint] OpenAI-style 表示仿 OpenAI API 形态的请求、流式事件或兼容服务；URL 是统一资源定位符；endpoint 是服务入口地址。兼容接口不代表底层 tokenizer、预算和取消语义自动一致。
  \item[TensorRT-LLM/SmoothQuant/LM head] TensorRT-LLM 是 NVIDIA 生态的大模型推理优化系统；SmoothQuant 是把 activation 和 weight 尺度迁移后再量化的一类方法；LM head 是把 hidden state 投影到词表 logits 的输出层。
  \item[PDF/QA] PDF 是固定版式文档格式；QA 在工程验收中常指质量保证，在问答应用中也可指 question answering。读到 QA 时要结合上下文判断是测试流程还是问答任务。
  \item[checkpoint/tool/chunk/bucket] checkpoint 是可恢复状态；tool 是模型可调用的外部工具；chunk 是文档或输入切片；bucket 是按范围、长度或哈希归类的桶。
  \item[scheduler/pipeline/target/epoch] scheduler 决定请求或 kernel 的执行顺序；pipeline 是多阶段处理链；target 是后端目标或质量/性能目标；epoch 是版本或轮次编号。
  \item[checksum/hash/contract/generation] checksum 是输出摘要；hash 是摘要或哈希值；contract 是输入输出和错误边界；generation 是同一对象的版本编号。
  \item[register/heap/stack/interrupt/owner] register 是寄存器；heap 是动态分配区域；stack 是调用栈；interrupt 是中断；owner 是当前拥有资源或写权限的一方。后端和服务代码里的所有权必须能被测试和日志说明。
  \item[roofline/throughput/bandwidth/coherence/shard/prefetch] roofline 是用算术强度估计性能上界的模型；throughput 是吞吐；bandwidth 是带宽；coherence 是缓存一致性；shard 是数据分片；prefetch 是提前取数据以隐藏延迟。
  \item[LayerNorm/ALU/TensorRT] LayerNorm 是按隐藏维做均值方差归一化的层；ALU 是算术逻辑单元；TensorRT 是 NVIDIA 推理优化生态，TensorRT-LLM 是其中面向大模型的系统。
  \item[sanitizer/worker] sanitizer 是编译器插桩诊断工具族，用于发现越界、未定义行为或数据竞争；worker 是执行请求、kernel 或后台任务的线程/执行单元。
  \item[MiniLLM/TinyTrace/TensorView] 这些是本册教学工程里的示例名：MiniLLM 表示最小本地推理线，TinyTrace 表示轻量事件追踪，TensorView 表示不拥有底层数据的张量视图。它们不是行业标准缩写，读到时按项目内对象理解。
  \item[Transformer 形状变量] \code{head\_dim} 是每个 attention head 的隐藏维宽度；\code{kv\_heads/kv\_head} 是 Key/Value head 数或其中一个 head；\code{query\_heads/q\_head} 是 Query head 数或其中一个 head；\code{n\_q/n\_kv} 常分别表示 query head 数和 KV head 数。看到这些变量时，先检查 MHA、GQA、MQA 三种形状关系。
  \item[KV cache 字段] \code{kv\_cache} 保存历史 K/V 张量；\code{prompt\_tokens} 是输入上下文 token 数；\code{max\_context} 和 \code{max\_new\_tokens} 分别限制上下文窗口和本次最多生成 token 数；\code{model\_position} 是模型看到的绝对或相对位置。它们共同决定能不能继续解码、会不会越界。
  \item[量化字段] \code{scale\_w} 和 \code{scale\_a} 常表示权重、激活的量化比例；\code{zero\_point} 表示整数零点；\code{group\_size} 表示共享一组量化参数的元素数量；\code{acc\_i32} 常表示 int8 乘加后用 32 位整数累加。量化 bug 往往来自这些字段的单位和广播方向不一致。
  \item[服务指标字段] \code{ttft\_ms} 是首 token 延迟毫秒数；\code{prefill\_tokens\_per\_s} 和 \code{decode\_tokens\_per\_s} 分别表示 prefill、decode 吞吐；\code{finish\_reason} 记录生成停止原因，例如到达 EOS、长度上限或取消请求。指标必须和请求大小一起解释。
  \item[后端实验名] \code{TILE\_N}、\code{TILE\_K}、\code{out\_tile}、\code{packed\_weight}、\code{tail\_policy} 这类名字是后端 schedule 和布局选择。它们不是模型层名，而是为了让内层循环更接近 SIMD、cache 和寄存器约束。
  \item[项目字段名] \code{request\_id}、\code{model\_id}、\code{layer\_id}、\code{block\_id}、\code{tokenizer\_hash}、\code{backend\_plan\_id} 这类 id/hash 字段用于重放和审计：id 定位对象，hash 确认内容版本，plan id 确认后端选择。模型工程不能只靠文件名猜版本。
  \item[数学边界短写] \code{K0/K1}、\code{V0/V1}、\code{W+1}、\code{T-1}、\code{G-1} 这类短写通常出现在 shape、窗口或循环边界推导中。字母本身没有固定全书含义，必须回到所在公式：K 常是 key 或归约维，V 常是 value，W 常是窗口，T 常是时间步或 token 位置，G 常是 group。
  \item[组合方向词] \code{GEMV/GEMM}、\code{QKV/MLP}、\code{ONNX/MLIR}、\code{TPOT/TTFT} 这类组合是在同一段比较几个已解释概念。斜杠表示并列，不表示一种新算法。
  \item[CMake 和转换] \code{cmake} 是构建系统命令；\code{static\_cast} 是 C++ 显式类型转换。第三册的工程验收章节用它们记录“如何构建”和“类型边界如何写清”，不是模型概念。
\end{description}
